\section{Related work}
\label{sec:related}

\paragraph{Metric depth and the camera.}
To predict metric depth from one image, a model has to deal with the focal length. Early work gave the
focal length to the network~\cite{he2018focal} or made the convolutions camera-aware~\cite{facil2019camconvs}.
FS-Depth~\cite{wei2023fsdepth} also augments focal lengths by cropping, with depth labels.
Today's models fall into three groups. Some \emph{take the camera}: Metric3D~\cite{yin2023metric3d,hu2024metric3dv2}
maps every image to a canonical camera and needs $K$ at test time, and Depth Any Camera~\cite{guo2025dac}
and ZeroDepth~\cite{guizilini2023zerodepth} are also conditioned on it. Some \emph{predict the camera}:
Depth Pro~\cite{bochkovskii2025depthpro} estimates the field of view,
UniDepth~\cite{piccinelli2024unidepth,piccinelli2025unidepthv2} predicts a ray per pixel, and
MoGe-2~\cite{wang2025moge2} returns a metric point map with intrinsics. The rest use \emph{pixels only},
such as the metric heads of Depth Anything~\cite{yang2024da,yang2024dav2} and ZoeDepth~\cite{bhat2023zoedepth}.
Training on many cameras~\cite{ma2026metricanything} is meant to make camera inference robust. Our audit
tests whether it is, and our three model groups follow these three families.

\paragraph{Depth cues.}
DepthCues~\cite{danier2025depthcues} probes which monocular cues large vision models pick up, among them
relative and familiar size. Our exponent $\alpha$ measures one cue in particular, apparent size, and our
cue intervention (\cref{sec:cues}) changes it directly while keeping everything else fixed.

\paragraph{Single-image calibration.}
Perspective Fields~\cite{jin2023perspectivefields} and related methods predict the camera from local
perspective cues, which is the kind of cue the zoom--squash ambiguity says a model must use. InFlux~\cite{liang2025influx}
is a benchmark for video cameras whose intrinsics change over time, for example by zooming, and shows that
most calibration methods struggle with it. We do not propose a calibration method. We ask whether depth
models respond correctly to the one camera change whose effect is known exactly, without any ground truth.

\paragraph{Crops and equivariance.}
Zhong \etal~\cite{zhong2023equivariance} are the closest to us. They show that supervised depth and normal
predictors are not equivariant to crop-and-resize: the prediction on a crop is not the crop of the
prediction. They fix it with test-time averaging and a self-consistency loss on unlabeled images. Their
view treats the output as a map that should simply be cropped, and the camera does not appear. For
\emph{metric} models the same requirement has a physical meaning. A crop is a zoom, the focal length must
grow by $s$, and a model that reads apparent size breaks the rule in a specific and predictable way. We
measure this with two exponents, explain it with the zoom--squash ambiguity, split it into its causes, and
test camera-predicting models that did not exist in 2023. Our fine-tuning (\cref{sec:eft}) follows the
self-consistency idea of~\cite{zhong2023equivariance} with a frozen teacher. We do not claim the loss as new,
and we compare against their stop-gradient target.
Kim and Lee~\cite{kim2026mismatch} study camera-conditioned metric models when a crop, resize or padding
changes the intrinsics but the pipeline still passes the old ones, and they propose a mismatch-aware camera
module. Our stale-$K$ run with Metric3D v2 is one case of this setting. We include it only to put
camera-conditioned models on the same scale as the other two groups.
Other work looks at crops in different tasks. PlaneDepth~\cite{wang2023planedepth} notes that resize-and-crop
augmentation breaks the plane geometry of self-supervised depth. KPE~\cite{prakash2024kpe} shows that crops
cause a perspective ambiguity for object-level 3D and feeds the crop position to the network. Multi-crop
human mesh recovery~\cite{nie2024multicrop} keeps the camera consistent between crops of one person.
NERO~\cite{zhao2023nero} visualises how models break equivariance along a transformation. A recent paper
makes depth models robust to camera roll~\cite{tang2026roll}. A roll does not change the intrinsics, but a
zoom does, so the right output under a zoom is not ``the same map, transformed'' but a specific change of
the metric prediction.

\paragraph{Label-free evaluation and adaptation.}
Self-Geometry~\cite{youn2026selfgeometry} adapts 3D foundation models at test time with multi-view
constraints, Re-Depth Anything~\cite{bhattarai2025redepth} refines depth by self-supervised re-lighting, and
test-time adaptation has been used for depth completion~\cite{park2024ttadepth}. A self-consistency score is
not an accuracy score, since a model can be consistently wrong. We therefore always report the label-free
exponents next to errors against ground truth, and we test when the first predicts the second
(\cref{sec:results}).
